% =========================================================
% Limits at Infinity — Notes
% Chapter: continuity
% Volume III · Analysis
% =========================================================

\subsection{Limits at Infinity}
\label{sec:limits-at-infinity}

\begin{tcolorbox}[title={Toolkit --- Limits at Infinity}]
\begin{itemize}
  \item $+\infty$ is adherent to $X$ iff $\neg\,\operatorname{BoundedSet}(X)$
    above ($\sup X=+\infty$).
  \item $\operatorname{LimitAtInfinity}(f,+\infty,L)$: replace
    $\operatorname{InputTolerance}(x,c,\delta)$ with $x>M$.
    Same $\operatorname{OutputTolerance}(f(x),L,\varepsilon)$.
  \item All standard limit laws carry over unchanged.
  \item Sequential characterisation: $\operatorname{LimitAtInfinity}
    (f,+\infty,L)\iff f(x_n)\to L$ for every $x_n\to+\infty$ in $X$.
\end{itemize}
\end{tcolorbox}

\begin{tcolorbox}[title={Definition --- Infinite Adherent Points}]
\begin{definition}[Infinite Adherent Points]
\label{def:infinite-adherent-points}
Let $X\subseteq\mathbb{R}$.
\begin{itemize}
  \item $+\infty$ \textbf{adherent to $X$}: $\forall M\in\mathbb{R},\;
    \exists x\in X:x>M$ $\iff\sup X=+\infty$.
  \item $-\infty$ \textbf{adherent to $X$}: $\forall M\in\mathbb{R},\;
    \exists x\in X:x<M$ $\iff\inf X=-\infty$.
\end{itemize}
$\operatorname{BoundedSet}(X)\iff$ neither $+\infty$ nor $-\infty$
is adherent to $X$.
\end{definition}
\end{tcolorbox}

\begin{remark*}[Standard quantified statement]
\[
  +\infty\text{ adherent to }X
  \;\iff\;
  \forall M\in\mathbb{R},\;\exists x\in X:x>M
  \;\iff\;
  \neg\,\operatorname{BoundedAbove}(X).
\]
\end{remark*}

\begin{remark*}[Definition predicate reading]
\[
  \operatorname{AdherentPoint}(+\infty,X)
  \;\colonequiv\;
  \forall M\in\mathbb{R},\;\exists x\in X:x>M.
\]
\end{remark*}

\begin{remark*}[Interpretation]
Both finite and infinite adherence are instances of adherence in
$\overline{\mathbb{R}}$: finite adherence uses balls; infinite
adherence uses rays.
\end{remark*}

\begin{tcolorbox}[title={Definition --- Limit at Infinity}]
\begin{definition}[Limit at Infinity]
\label{def:limit-at-infinity}
Let $X\subseteq\mathbb{R}$ with $\operatorname{AdherentPoint}(+\infty,X)$,
$f:X\to\mathbb{R}$.
$\operatorname{LimitAtInfinity}(f,+\infty,L)$ if
\[
  \forall\varepsilon>0,\;
  \exists M\in\mathbb{R},\;
  \forall x\in X:\;
  x>M \;\Rightarrow\; |f(x)-L|<\varepsilon.
\]
\end{definition}
\end{tcolorbox}

\begin{remark*}[Standard quantified statement]
\[
  \operatorname{LimitAtInfinity}(f,+\infty,L)
  \;\iff\;
  \forall\varepsilon>0,\;\exists M\in\mathbb{R},\;\forall x\in X:\;
  x>M \Rightarrow |f(x)-L|<\varepsilon.
\]
\end{remark*}

\begin{remark*}[Negated quantified statement]
\[
  \neg\,\operatorname{LimitAtInfinity}(f,+\infty,L)
  \;\iff\;
  \exists\varepsilon_0>0,\;\forall M\in\mathbb{R},\;\exists x\in X:\;
  x>M \;\wedge\; |f(x)-L|\geq\varepsilon_0.
\]
\end{remark*}

\begin{remark*}[Comparison with finite limits]
$\operatorname{FunctionLimit}(f,c,L)$ uses
$\operatorname{InputTolerance}(x,c,\delta)$.
$\operatorname{LimitAtInfinity}(f,+\infty,L)$ replaces this with
$x>M$. The output condition is identical in both. All limit laws
carry over with the same proofs.
\end{remark*}

\begin{remark*}[Interpretation]
The threshold $M$ plays the role of $c-\delta$ in the finite limit:
everything to the right of $M$ in $X$ must lie within
$\varepsilon$ of $L$.
\end{remark*}
